\section{Error Correction: The Steane Code}\label{sec:error-correction}

This section first provides a concise overview of quantum error correction principles to establish context for readers less familiar with these techniques. Following this brief introduction, we detail our specific implementation of the Steane code within the qstack framework. Our complete implementation will be provided as supplementary material and consists of a single Python file of about 200 lines of well-formatted code, with most of it dedicated to describing the output circuits.

\subsection{Quantum Error Correction}

As with classical counterparts, quantum codes work by encoding one qubit of information into multiple to achieve redundancy and resist individual errors. However, several uniquely quantum obstacles complicate this approach:

\begin{enumerate}
\item Due to the no-cloning theorem, we can’t copy quantum information into multiple qubits.
\item We can’t directly measure the state of a qubit without disturbing it.
\item Quantum errors aren’t just bit flips — they can include phase flips, or both simultaneously.
\end{enumerate}

To work around these constraints, redundancy is introduced not by copying, but by entanglement. For example, a single-qubit state $\alpha \ket{0} + \beta \ket{1}$ can be encoded as $\alpha \ket{000} + \beta \ket{111}$. Instead of inspecting individual qubit states, which would collapse the superposition, we perform parity checks across groups of qubits. These can be implemented indirectly using ancilla qubits to determine whether two or more qubits are in the same state (that is, have even parity), without revealing what that state actually is.

This use of entanglement makes quantum codes remarkably sensitive: a single-qubit error can disturb the entire encoded state. Consider, for instance, the three-qubit state $\alpha \ket{000} + \beta \ket{111}$. If a bit-flip error (Pauli $X$) occurs on the first qubit, the state becomes $\alpha \ket{100} + \beta \ket{011}$, which is still entangled but no longer lies in the original code space.

These parity-check operations are called stabilizers. A stabilizer is a quantum operator that leaves valid encoded states invariant, that is, the state is a $+1$ eigenstate of the stabilizer. Stabilizers composed of Pauli-$Z$ operators are used to detect bit flip (Pauli-$X$) errors, and are called $Z$ stabilizers. Similarly, stabilizers made from Pauli-$X$ operators are used to detect phase flip (Pauli-$Z$) errors, and are called $X$ stabilizers. More generally, some codes use mixed stabilizers with both $X$ and $Z$ terms to detect a broader class of errors.

A quantum error-correcting code is typically defined by a set of stabilizer operators. A quantum state is considered valid (i.e., free of detectable errors) if it is a $+1$ eigenstate of every stabilizer in the set. Measuring stabilizers doesn’t collapse the encoded state, it only tells us whether the parity constraints still hold. The result of all stabilizer measurements is collected into a bitstring called the syndrome. This syndrome doesn’t identify which specific qubit failed, but it narrows the possibilities. The syndrome decoder is the map from each syndrome to the most likely error (MLE), which can then be corrected, and the procedure that performs all measurements is the syndrome extraction routine.

In summary, the process of quantum error correction involves:

\begin{enumerate}
\item \textbf{Encoding}: spread one logical qubit across multiple physical ones using entanglement.
\item \textbf{Syndrome extraction}: measure stabilizers to detect signs of corruption without collapsing the logical state.
\item \textbf{Syndrome decoding}: analyze the syndrome to determine the most probable error.
\item \textbf{Correction}: apply the inverse operation to return to the original code space.
\end{enumerate}

As with classical codes, quantum codes are characterized by three parameters:

\begin{itemize}
\item $n$ = number of physical qubits used by the code
\item $k$ = number of logical qubits encoded
\item $d$ = code distance, the minimum number of physical qubits that must be altered to implement a nontrivial logical operation or transform one valid encoded state into another (i.e., smallest weight of a logical operator not in the stabilizer group)
\end{itemize}

When describing the characteristics of a code, these parameters are denoted as:

$$
[[n, k, d]]
$$


To demonstrate quantum error correction within our framework, we implemented the Steane code, a \([[7, 1, 3]]\) CSS (Calderbank-Shor-Steane) code. The Steane code encodes 1 logical qubit into 7 physical qubits and can correct any single-qubit phase or bit flip error. Based on the classical Hamming code, it features a symmetric structure where X and Z stabilizers mirror each other, enabling transversal implementations of logical $H$ and $CX$ operations ---a valuable property for fault-tolerant quantum computing.

The code is defined by the following stabilizer generators:~
\[
\begin{aligned}
\text{Z stabilizers: } &[Z, Z, I, Z, Z, I, I] \\
&[Z, I, Z, Z, I, Z, I] \\
&[I, Z, Z, Z, I, I, Z] \\
\text{X stabilizers: } &[X, X, I, X, X, I, I] \\
&[X, I, X, X, I, X, I] \\
&[I, X, X, X, I, I, X] \\
\end{aligned}
\]

The logical operators that act on the encoded qubit are:~
\[
X_L = [X, X, X, X, X, X, X], \quad
Z_L = [Z, Z, Z, Z, Z, Z, Z]
\]

\subsection{Compilation Process}

Our compiler transforms quantum programs to incorporate error correction transparently, following the compositional compilation approach introduced in earlier sections. When a kernel in the source program allocates a logical qubit, the compiler applies several transformations:

\begin{enumerate}
\item \textbf{Qubit expansion}: A single logical qubit is expanded into 7 physical qubits
\item \textbf{State preparation}: The compiler injects a preparation circuit that initializes the quantum system to the $+1$ eigenspace of all Z stabilizers, representing the logical zero state $\ket{0}_L$
\item \textbf{Callback injection}: The compiler wraps the original callback, by invoking a \texttt{decode} method that processes measurement outcomes to calculate a single logical measurement.
\end{enumerate}

\begin{figure}[t]
\begin{minipage}{0.48\textwidth}
\begin{minted}[fontsize=\scriptsize]{ruby}
allocate q1_0 q1_1 q1_2 q1_3 q1_4 q1_5 q1_6:
  h q1_4
  h q1_5
  h q1_6
  cx q1_4 q1_0
  cx q1_4 q1_1
  cx q1_4 q1_3
  cx q1_5 q1_0
  cx q1_5 q1_2
  cx q1_5 q1_3
  cx q1_6 q1_1
  cx q1_6 q1_2
  cx q1_6 q1_3
measure
\end{minted}
\centering{(a) State preparation circuit}
\end{minipage}
\hfill
\begin{minipage}{0.48\textwidth}
\begin{minted}[fontsize=\scriptsize]{ruby}
allocate q1_z_0 allocate q1_z_1 q1_z_2:
  cx q1_0 q1_z_0
  cx q1_1 q1_z_0
  cx q1_3 q1_z_0
  cx q1_4 q1_z_0
  cx q1_0 q1_z_1
  cx q1_2 q1_z_1
  cx q1_3 q1_z_1
  cx q1_5 q1_z_1
  cx q1_1 q1_z_2
  cx q1_2 q1_z_2
  cx q1_3 q1_z_2
  cx q1_6 q1_z_2
measure correct_z(qubit=q1)
\end{minted}
\centering{(b) Z-stabilizer syndrome extraction}
\end{minipage}
\caption{Quantum circuit components of the Steane code implementation: (a) The state preparation circuit that initializes the 7 physical qubits into the encoded $\ket{0}_L$ state, and (b) Z-stabilizer syndrome extraction circuit that uses ancilla qubits to detect bit-flip errors without disturbing the logical state.}
\label{fig:steane-code-circuits}
\end{figure}

When applied to a simple allocation and measurement of a qubit, our compiler generates a circuit that prepares the encoded state across 7 physical qubits (shown in Figure~\ref{fig:steane-code-circuits}a). Even empty kernels with no logical instructions undergo this state preparation to ensure proper initialization of the encoded logical state. The approach resembles the \texttt{vote} callback from the majority vote example (Section~\ref{sec:majority-vote}), but employs more sophisticated error detection logic tailored to the Steane code's structure.

A critical feature of our approach is that the compiled error-corrected program preserves the semantics of the original. Although the compilation process introduces significant complexity in the quantum circuit, the semantics remains unchanged—a single logical measurement outcome is pushed onto the measurement stack, enabling any classical postprocessing code from the original program to behave identically.

During compilation, each instruction in the original source \textbf{is replaced by its encoded counterpart}. For many operations in the Steane code, these mappings are simple transversal operations: applied identically to each physical qubit.

The compiler strategically inserts \textbf{syndrome extraction} and correction routines \textit{between} logical instructions. These routines identify and correct errors that may have occurred during computation.

For syndrome extraction, our implementation uses one ancilla qubit per stabilizer measurement, a straightforward approach that provides clear pedagogical value. Figure~\ref{fig:steane-code-circuits}b shows an example Z-stabilizer syndrome extraction circuit. These implementations can be trivially replaced with fault-tolerant versions. The circuit follows this pattern:

\begin{enumerate}
\item Initialize ancilla qubits in the $\ket{0}$ state for measuring Z stabilizers (or $\ket{+}$ for X stabilizers)
\item Entangle each ancilla with the relevant data qubits using controlled-Pauli gates according to the stabilizer definition
\item Measure the ancilla qubits to extract syndrome bits without disturbing the logical state
\end{enumerate}

Each syndrome extraction measurement triggers a compiler-injected callback that:
\begin{itemize}
\item Processes the syndrome bits to identify error locations
\item Returns a kernel with appropriate correction operations when errors are detected
\item Returns nothing when no corrections are needed
\end{itemize}

After each logical qubit measurement, the \texttt{decode} method is inserted before invoking the original callback. This method implements a complete error detection and correction workflow, as illustrated in Figure~\ref{fig:steane-code-decode}. It calculates syndrome bits by computing parity values according to the stabilizers, identifies the most likely error using the syndrome, and applies the necessary correction. Our implementation uses a static lookup table for syndrome-to-error mapping, which is effective for small codes like Steane, though larger codes would benefit from more sophisticated decoders such as minimum-weight perfect matching or neural networks~\cite{higgott2022pymatching, yao2024belief}.

\begin{figure}[t]
\begin{minted}[fontsize=\scriptsize]{python}
def decode(m1, m2, m3, m4, m5, m6, m7):
    outcome = np.array([m1, m2, m3, m4, m5, m6, m7])
    check1 = int(np.dot(outcome, np.array([1, 1, 0, 1, 1, 0, 0])) % 2)
    check2 = int(np.dot(outcome, np.array([1, 0, 1, 1, 0, 1, 0])) % 2)
    check3 = int(np.dot(outcome, np.array([0, 1, 1, 1, 0, 0, 1])) % 2)
    syndrome = (check1, check2, check3)
    fault = syndrome_table.get(syndrome)

    if fault is not None:
        outcome[fault] = (outcome[fault] + 1) % 2

    return int(np.sum(outcome) % 2)
\end{minted}
\caption{The decoder for the Steane code implementation processes measurement outcomes, calculates syndromes, and applies corrections based on a lookup table before determining the logical measurement result.}
\label{fig:steane-code-decode}
\end{figure}

The complete compilation transformation elegantly integrates physical-level error correction into higher-level quantum programming without exposing implementation details to the programmer.

What makes our approach particularly powerful is its composability with other transformations. The Steane code implementation can be composed with other compilation passes, enabling a separation of concerns where error correction becomes just another layer in the compilation pipeline. This aligns with modern software engineering practices where modularity and layered abstractions are essential for managing complexity.
